\chapter*{Sumber dan Atribusi}
\addcontentsline{toc}{chapter}{Sumber dan Atribusi}
\markboth{Sumber dan Atribusi}{Sumber dan Atribusi}

Buku ini mengadaptasi materi dari beberapa buku teks dan catatan kuliah
sumber terbuka. Setiap sumber yang dicantumkan di sini menggunakan lisensi
yang kompatibel dengan lisensi Creative Commons Attribution-ShareAlike~4.0
International (CC~BY-SA~4.0) yang menjadi dasar distribusi buku ini. Setiap
entri mencatat penulis asli, URL sumber, lisensi asli, dan cara materi tersebut
digunakan. Modifikasi terhadap seluruh materi yang diadaptasi dilakukan oleh
Robert Hildebrand dan kontributor lain buku ini.

%-------------------------------------------------------------------
\section*{Buku Teks dan Catatan Sumber Terbuka yang Diadaptasi}
%-------------------------------------------------------------------

\paragraph{A First Course in Linear Algebra}\hspace{0pt}\\[-0.6em]
\begin{description}
  \item[Penulis:] Ken Kuttler, dengan adaptasi oleh Lyryx Learning
  \item[Sumber:] \url{https://lyryx.com/first-course-linear-algebra/}
  \item[Lisensi:] Creative Commons Attribution 4.0 International (CC~BY~4.0)
  \item[Salinan lokal:] \texttt{external-sources/aFirstCourseLinearAlgebra/}
  \item[Digunakan dalam:] Lampiran tentang Aljabar Linear (vektor, hasil
    kali titik, hasil kali silang, proyeksi, Gram--Schmidt, bidang, dan
    garis parametrik). Buku ini juga menggunakan kembali templat
    pemformatan bab dan teorema milik Lyryx.
  \item[Kompatibilitas:] Materi CC~BY~4.0 dapat dimasukkan ke dalam karya
    CC~BY-SA~4.0 dengan atribusi. Materi yang diadaptasi di sini
    dilisensikan ulang berdasarkan CC~BY-SA~4.0 sebagai bagian dari buku ini.
\end{description}

\paragraph{Foundations of Applied Mathematics: Volume 2 -- Algorithms,
  Approximation, and Optimization}\hspace{0pt}\\[-0.6em]
\begin{description}
  \item[Penulis:] Jeffrey Humpherys, Tyler Jarvis, dan para kontributor di
    Brigham Young University (daftar lengkap kontributor tersedia di repositori
    sumber berikut; ucapan terima kasih ringkas dicantumkan di bagian akhir)
  \item[Sumber:] \url{https://github.com/Foundations-of-Applied-Mathematics}
  \item[Lisensi:] Creative Commons Attribution 3.0 United States License
    (CC~BY~3.0~US)
  \item[Salinan lokal:] \texttt{external-sources/foundationsAppliedMathematicsLabs/}
  \item[Digunakan dalam:] Bagian-bagian terpilih, templat pemformatan kode,
    dan rujukan ke kode laboratorium terkait di seluruh buku.
  \item[Kompatibilitas:] CC~BY~3.0 mengizinkan adaptasi berdasarkan lisensi
    apa pun yang kompatibel; lisensi pengadaptasi di sini adalah CC~BY-SA~4.0.
    Atribusi dipertahankan untuk seluruh materi yang diadaptasi.
\end{description}

\paragraph{Linear Inequalities and Linear Programming}\hspace{0pt}\\[-0.6em]
\begin{description}
  \item[Penulis:] Kevin Cheung, Carleton University
  \item[Sumber:] \url{https://github.com/dataopt/lineqlpbook}
  \item[Lisensi:] Creative Commons Attribution-ShareAlike 4.0
    International (CC~BY-SA~4.0)
  \item[Salinan lokal:] \texttt{external-sources/lineqlpbook/}
  \item[Digunakan dalam:] Bagian-bagian penting dari Bagian~I (Pemrograman
    Linear): notasi, solusi fisibel dasar, Teorema Fundamental LP, lema
    Farkas, kelonggaran komplementer, dualitas, metode eliminasi
    Fourier--Motzkin, dan contoh grafis. Lisensi didokumentasikan melalui
    kepala berkas pada setiap berkas sumber \texttt{.tex}.
  \item[Kompatibilitas:] Konten sumber CC~BY-SA~4.0 secara langsung
    kompatibel dengan adaptasi CC~BY-SA~4.0; tidak ada transisi versi.
\end{description}

\paragraph{Catatan Optimisasi (Catatan Pemrograman Linear)}\hspace{0pt}\\[-0.6em]
\begin{description}
  \item[Penulis:] Douglas Bish, Virginia Tech
  \item[Lisensi:] Creative Commons Attribution-ShareAlike 4.0
    International (CC~BY-SA~4.0)
  \item[Salinan lokal:] \texttt{external-sources/Opt\_Notes/}
  \item[Digunakan dalam:] Bagian-bagian penting dari Bagian~I (Pemrograman
    Linear), termasuk pemodelan, uraian metode simpleks, dan analisis
    sensitivitas. Lisensi didokumentasikan dalam
    \texttt{external-sources/Opt\_Notes/README.md}.
  \item[Kompatibilitas:] Konten sumber CC~BY-SA~4.0 secara langsung
    kompatibel dengan adaptasi CC~BY-SA~4.0.
\end{description}

\paragraph{Math in Society, bab Teori Graf}\hspace{0pt}\\[-0.6em]
\begin{description}
  \item[Penulis:] David Lippman, Pierce College Ft.\ Steilacoom
  \item[Sumber:] \url{http://www.opentextbookstore.com/mathinsociety/}
  \item[Lisensi:] Creative Commons Attribution-ShareAlike 3.0 United
    States (CC~BY-SA~3.0)
  \item[Digunakan dalam:] Bab Algoritme Graf (Buku 1): contoh motivasi
    pemeriksa halaman rumput dan K\"onigsberg, definisi dasar graf, contoh
    penyelesaian dengan algoritme Dijkstra dan pohon rentang berbiaya minimum, serta banyak
    latihan bab, semuanya dengan suntingan dan penambahan.
  \item[Foto sumber:] Foto udara pembangunan perumahan Pleasant View karya
    Sam Beebe, \url{https://www.flickr.com/photos/sbeebe/2850476641/},
    digunakan berdasarkan Creative Commons Attribution~2.0 Generic
    (CC~BY~2.0). David Lippman menambahkan penyajian dan anotasi graf pada
    gambar turunan dalam \emph{Math in Society}.
  \item[Kompatibilitas:] Materi CC~BY-SA~3.0 dapat diadaptasi dan
    didistribusikan ulang berdasarkan CC~BY-SA~4.0, yang dicantumkan oleh
    Creative Commons sebagai versi lisensi lebih baru yang kompatibel untuk
    adaptasi. Atribusi dipertahankan di sini dan dalam berkas sumber bab.
\end{description}

\paragraph{Applied Finite Mathematics, Bab 1: Persamaan Linear}\hspace{0pt}\\[-0.6em]
\begin{description}
  \item[Penulis:] Rupinder Sekhon dan Roberta Bloom
  \item[Sumber:] \url{https://math.libretexts.org/Bookshelves/Applied_Mathematics/Applied_Finite_Mathematics_(Sekhon_and_Bloom)}
  \item[Lisensi:] Creative Commons Attribution 4.0 International
    (CC~BY~4.0), melalui LibreTexts
  \item[Salinan lokal:] \texttt{external-sources/applied-finite-mathematics/}
    (dengan berkas LICENSE)
  \item[Digunakan dalam:] Lampiran A (Persamaan Linear), termasuk teks dan
    gambar yang diadaptasi dari Bab~1 LibreTexts.
  \item[Kompatibilitas:] Materi CC~BY~4.0 dapat dimasukkan ke dalam karya
    CC~BY-SA~4.0 dengan atribusi. Paragraf atribusi yang terlihat
    dicantumkan pada awal bab lampiran.
\end{description}

%-------------------------------------------------------------------
\section*{Atribusi Gambar}
%-------------------------------------------------------------------

\paragraph{Jembatan K\"onigsberg.}
Ukiran bersejarah tujuh jembatan (digunakan dalam bab Teori Graf) berada
dalam domain publik. Varian gambar garis yang disertakan di sini dibuat
oleh Bogdan Giusca (pengguna Wikipedia Bogdang) dan dilisensikan berdasarkan
Creative Commons Attribution-ShareAlike~3.0 (CC~BY-SA~3.0).
Sumber: \url{https://commons.wikimedia.org/wiki/File:Konigsberg_bridges.png}.

\paragraph{Ilustrasi masalah ransel.}
Ilustrasi ransel dan empat kotak berasal dari ``Knapsack.svg'' karya pengguna
Wikimedia Commons Dake, dengan suntingan berikutnya oleh Keenan Pepper, dan
digunakan berdasarkan Creative Commons Attribution-ShareAlike~2.5 Generic
(CC~BY-SA~2.5). Sumber dan riwayat berkas:
\url{https://commons.wikimedia.org/wiki/File:Knapsack.svg}.

\paragraph{Gambar asli dan pengecualian pihak ketiga.}
Plot dan diagram yang pada metadata per berkas mencantumkan Robert Hildebrand
atau kontributor buku sebagai pembuat merupakan karya asli mereka. Keempat
gambar optimisasi multiobjektif dalam Bab~13 merupakan karya asli Robert
Hildebrand. Gambar yang berasal dari \emph{Foundations of Applied
Mathematics}, \emph{Math in Society}, Wikimedia Commons, atau sumber pihak
ketiga lain mempertahankan atribusi dan lisensi khususnya pada metadata
\texttt{00\_METADATA.bib}; lisensi umum buku tidak menggantikan lisensi khusus
tersebut.

%-------------------------------------------------------------------
\section*{Ucapan Terima Kasih Lainnya}
%-------------------------------------------------------------------

Kami berterima kasih kepada Laurent Poirrier dan Diego Moran atas kontribusi
catatan mengenai pemrograman linear dan bilangan bulat, serta kepada Jamie
Fravel atas penyuntingan dan kontribusi bab serta kode contoh. Kami juga
berterima kasih kepada Irma Adams karena telah meneliti dan menulis kotak
samping biografi ``sorotan perintis'' yang muncul di seluruh buku (termasuk
profil George Dantzig dan para pendiri bidang ini lainnya). Joshua Emmanuel
memberikan arahan awal mengenai penyajian metode grafis selama pengembangan
mata kuliah; teks saat ini mengenai metode grafis merupakan karya asli buku ini.

%-------------------------------------------------------------------
\section*{Ringkasan Kompatibilitas Lisensi}
%-------------------------------------------------------------------

\begin{center}
\renewcommand{\arraystretch}{1.25}
\begin{tabular}{l l l}
\toprule
\textbf{Lisensi sumber} & \textbf{Kompatibilitas} & \textbf{Mekanisme} \\
\midrule
CC~BY-SA~4.0 & Langsung & Lisensi sama; atribusi dipertahankan \\
CC~BY~4.0    & Langsung & Diadaptasi ke CC~BY-SA~4.0 dengan atribusi \\
CC~BY~3.0~US & Tidak langsung & Ketentuan lisensi pengadaptasi mengizinkan CC~BY-SA~4.0 \\
CC~BY~2.0 & Langsung & Atribusi dan pranala lisensi dipertahankan \\
CC~BY-SA~2.5 & Berbagi serupa & Adaptasi didistribusikan berdasarkan CC~BY-SA~4.0 \\
CC~BY-SA~3.0~US & Berbagi serupa & Adaptasi didistribusikan berdasarkan CC~BY-SA~4.0 \\
Domain Publik & Langsung & Tanpa pembatasan; atribusi diberikan \\
\bottomrule
\end{tabular}
\end{center}

\noindent Teks lengkap lisensi CC~BY-SA~4.0 yang menjadi dasar distribusi
buku ini tersedia di \url{https://creativecommons.org/licenses/by-sa/4.0/}
dan dalam berkas \texttt{LICENSE-Content} di repositori sumber buku.
Kode yang dirilis bersama buku ini dilisensikan secara terpisah berdasarkan
\textbf{MIT License}, dengan teks lisensi dalam \texttt{LICENSE-Code}.

\cleardoublepage
